% Circuit symbols used by the direct-current exercises.
\tikzset{
  em resistor/.style={
    draw,fill=white,minimum width=.64cm,minimum height=.19cm,inner sep=0pt,
  },
}
\newcommand{\EMfourthDiagram}[1]{%
  \adjustbox{max totalsize={.87\linewidth}{.24\textheight}}{%
    \input{tikzfigs/chapter4/fig_#1.tex}%
  }%
}
\newcommand{\EMdot}[1]{\fill (#1) circle[radius=1.25pt];}
\newcommand{\EMterm}[3]{%
  \draw[fill=white] (#1) circle[radius=1.65pt];
  \if\relax\detokenize{#3}\relax\else
    \node[#2,outer sep=1.65pt] at (#1) {$#3$};
  \fi
}
\newcommand{\EMres}[4][above=4pt]{%
  \draw (#2)--node[midway,sloped,em resistor]{} (#3);
  \if\relax\detokenize{#4}\relax\else
    \path (#2)--node[midway,sloped,#1,inner sep=1pt] {$#4$} (#3);
  \fi
}
\newcommand{\EMmeter}[2]{%
  \draw[fill=white] (#1) circle[radius=.31] node[inner sep=0pt] {$#2$};
}
% The long plate is the positive terminal before rotation.
\newcommand{\EMcell}[2][0]{%
  \begin{scope}[shift={(#2)},rotate=#1]
    \draw[white,line width=2pt] (-.11,0)--(.11,0);
    \draw (-.11,-.26)--(-.11,.26) (.11,-.17)--(.11,.17);
  \end{scope}%
}
\newcommand{\EMswitch}[2][0]{%
  \begin{scope}[shift={(#2)},rotate=#1]
    \draw[white,line width=2pt] (-.25,0)--(.25,0);
    \draw (-.25,0)--(.17,.26);
    \EMterm{-.25,0}{}{}
    \EMterm{.25,0}{}{}
  \end{scope}%
}
% Four arms of a rectangular bridge. The center branch is drawn at the call site.
\newcommand{\EMbridge}[4]{%
  \draw (-.64,0)--(0,0) (4,0)--(4.64,0)
    (0,-.88)--(0,.88) (4,-.88)--(4,.88);
  \EMres{0,.88}{2,.88}{#1}
  \EMres{2,.88}{4,.88}{#2}
  \EMres[below=4pt]{0,-.88}{2,-.88}{#3}
  \EMres[below=4pt]{2,-.88}{4,-.88}{#4}
  \EMdot{0,0}
  \EMdot{4,0}
}
% Hexagonal resistor network with terminals on its horizontal diagonal.
\newcommand{\EMhexagon}{%
  \foreach \vertex/\ang in {1/180,2/120,3/60,4/0,5/-60,6/-120}
    \coordinate (hex\vertex) at (\ang:2.15);
  \foreach \a/\b in {1/2,2/3,3/4,1/6,6/5,5/4} {
    \EMres{hex\a}{hex\b}{r}
  }
  \draw (-2.78,0)--(hex1) (hex4)--(2.78,0);
  \EMterm{-2.78,0}{left}{a}
  \EMterm{2.78,0}{right}{b}
  \foreach \vertex in {1,...,6} { \EMdot{hex\vertex} }
}
